\subsection{函子性质}

在本号中，所有代数均假设为结合的且含单位元的，并且每个代数同态均假设为含单位元的。

命题 9. 设 A、B 为两个分次 K-代数，E 为 (A, A)-双模，F 为分次 (B, B)-双模；设 \(\rho: \mathbf{A} \to \mathbf{B}\) 是一个分次代数同态，并且 \(\theta: \mathbf{E} \to \mathbf{F}\) 是 A-双模（相对于 \(\rho\) ）的一个分次 A-同态，次数为 0。则：

\begin{enumerate}
\def\labelenumi{(\roman{enumi})}
\item
  对每个 \(\varepsilon\) -导子 \(d':\mathbf{B}\to\mathbf{F}\) ， \(d'\circ\rho:\mathbf{A}\to\rho^{*}(\mathbf{F})\) 是一个同次数的 \(\varepsilon\) -导子。
\item
  对每个 \(\varepsilon\) -导子 \(d: \mathbf{A} \to \mathbf{E}\) ， \(\theta \circ d: \mathbf{A} \to \rho^{*}(\mathbf{F})\) 是一个同次数的 \(\varepsilon\) -导子。
\end{enumerate}

这两个断言直接从这些关系中得出。

\[
\begin{array}{l} d ^ {\prime} (\rho (x y)) = d ^ {\prime} (\rho (x) \rho (y)) = d ^ {\prime} (\rho (x)) \rho (y) + \varepsilon_ {\delta ^ {\prime}, \xi} \rho (x) d ^ {\prime} (\rho (y)) \\ \theta (d (x y)) = \theta ((d x) y + \varepsilon_ {\delta, \xi} x (d y)) = \theta (d x) \rho (y) + \varepsilon_ {\delta, \xi} \rho (x) \theta (d y) \end{array}
\]

其中 \(x \in \mathbf{A}\) 是次数为 \(\xi\) 的齐次元素， \(y \in \mathbf{A}\) ，而 \(\delta\) 与 \(\delta'\) 分别表示 \(d\) 与 \(d'\) 的次数。

推论。设 S 是代数 A 的一个生成系。为使 \(d' \circ \rho = \theta \circ d\) ，必要且充分条件是 \(d'(\rho(x)) = \theta(d(x))\) 对所有 \(x \in S\) 成立。

这是命题 9 以及第 5 号命题 3 的推论的直接推论。

在命题 9 的条件下，我们知道 B（借助 \(\rho\) ）具有一个 (A, A)-双模结构（II，§ 1，第 14 号，例 1）。

命题 10. 在命题 9 的条件下，要使 \(\varepsilon\) -导子 \(d': \mathbf{B} \to \mathbf{F}\) 对 B 与 \(\rho_*(\mathbf{F})\) 上的左（相应地，右）A-模结构是 A-线性的，必要且充分条件是 \(d'\) 在 B 的子代数 \(\rho(\mathbf{A})\) 上为零。

我们对左 A-模结构进行证明。对 \(a \in \mathbf{A}\) ， \(b \in \mathbf{B}\) ，

\[
d ^ {\prime} (\rho (a) b) = d ^ {\prime} (\rho (a)) b + \rho (a) d ^ {\prime} b
\]

于是，若 \(d' \circ \rho = 0\) ，则 \(d'\) 对 B 与 \(\rho^*(\mathbf{F})\) 上的左 A-模结构是线性的。反之，若确实如此，则特别地

\[
d ^ {\prime} (\rho (a)) = d ^ {\prime} (\rho (a). 1) = \rho (a) d ^ {\prime} (1) = 0
\]

（第 5 号，命题 3）。

特别地，设 \(D_{K}(B,F)\) 表示 B 到 F 的导子组成的 K-模（第 2 号）；其中 A-线性的那些导子，换句话说在 \(\rho(A)\) 上为零的那些导子，构成 \(D_{K}(B,F)\) 的一个子 K-模，记作 \(D_{A,\rho}(B,F)\) 或简记作 \(D_{A}(B,F)\) （显然 \(D_{K}(B,F)=D_{K,\phi}(B,F)\) ，其中 \(\phi:K\to B\) 是定义 B 上 K-代数结构的同态）。

现在设 A、B、C 为三个分次 K-代数， \(\rho:A\to B\) ， \(\sigma:B\to C\) 为两个分次代数同态，G 为一个分次 (C, C)-双模；如果 \(D_{A}(B,G)\) , \(D_{B}(C,G)\) 和 \(D_{A}(C,G)\) 表示相应的 K-模 \(D_{A,\rho}(B,\sigma_{*}(G))\) , \(D_{B,\sigma}(C,G)\) 和 \(D_{A,\sigma_{o}\rho}(C,G)\) ，则 \(D_{B}(C,G)\) 显然是 \(D_{A}(C,G)\) 的一个子 K-模，因为 \(\sigma(\rho(A))\subset\sigma(B)\) 。

命题 11. 在上述条件下，存在 K-同态的一个正合序列

\[
0 \rightarrow \mathrm{D} _ {\mathrm{B}} (\mathrm{C}, \mathrm{G}) \xrightarrow {u} \mathrm{D} _ {\mathrm{A}} (\mathrm{C}, \mathrm{G}) \xrightarrow {v} \mathrm{D} _ {\mathrm{A}} (\mathrm{B}, \mathrm{G})\tag{27}
\]

其中 u 是典范单射，v 是同态 \(d \mapsto d \circ \sigma\) （命题 9）。v 的核是这样的导子 \(d: C \to G\) 组成的集合：\(d(\sigma(b)) = 0\) 对所有 \(b \in B\) 成立；它正是 u 的像。

